%%%%%%%%%%%%%%%%%%%%%%%%% Abbreviations %%%%%%%%%%%%%%%%%%%%%%%%%%%%%%%%%%%%%%%%%
\chapter*{List of abbreviations}
\addcontentsline{toc}{chapter}{List of abbreviations}
\fancyhead[RE,LO]{List of abbreviations}
\fancyhead[RO,LE]{\thepage}

%%%%%%%%%%%%%%%%%%%%%%%%% How to add abbreviations %%%%%%%%%%%%%%%%%%%%%%%%%%%%%
% For each abbreviation you are writing, add into the acronym environment, e.g.:
% \acro{EA}{Expanded abbreviation}.
% If you have opted to have integrated acronyms with option 1 in the preamble,
% you can call this with \ac{EA}.
% If you need to add further information, use the \acroextra command as per this
% example from Jonathan Coney:
% \acro{BR}{British Rail\acroextra{, the state-owned operator of railways in Britain during the second half of the 20th century}}


\begin{acronym}[HALO-(AC)$^{3}$]\itemsep4pt  
% ^ in the square brackets above, add your longest abbreviatoin. This dictates
% the amount of spacing in the "table".
% Unfortunately, you need to manually keep track of the alphabetical order
% but hopefully, you know the alphabet.
% Don't forget this is a list of abbreviations, so abbreviated words must
% be included as well as initialisms.
\acro{CAO}{cold-air outbreak}
\acro{HALO-(AC)$^{3}$}{Project on Arctic Amplification Climate Relevant Atmospheric and Surface Processes and Feedback Mechanisms}
\acro{lat.}{latitude}
\end{acronym}
